\hypersetup{pageanchor=false}

\begin{titlepage}
    \centering
    {\large Universidade Presbiteriana Mackenzie \par}
    \vspace{0.4cm}
    {\large Banco de Dados – Projeto Aplicado II \par}
    \vspace{2.0cm}
    {\Large \textbf{Etapa 2} \par}
    \vspace{0.4cm}
    {\large \textbf{Definição do Método Analítico} \par}
    \vspace{1.4cm}
    {\Large \textbf{Classificação e Monitoramento Textual de Boletins Epidemiológicos sobre Dengue no Brasil} \par}
    \vfill

    \begin{flushleft}
    \textbf{Grupo:} Ana Clara Silva de Souza; Cid Wallace Araujo de Oliveira; Eduardo Machado Silva; Frederico Ripamonte Borges \\
    \textbf{Disciplina:} Projeto Aplicado II \\
    \textbf{Professor:} Anderson Adaime de Borba \\
    \textbf{Semestre:} 2026.01
    \end{flushleft}

    \vfill
    {\large \today \par}
\end{titlepage}

\clearpage
\hypersetup{pageanchor=true}
\setcounter{page}{1}

\section*{Resumo Executivo}

Esta segunda etapa do projeto apresenta a definição do método analítico que será utilizado para a classificação textual supervisionada de boletins epidemiológicos relacionados a arboviroses no Brasil. Em coerência com a proposta da Etapa 1, a pesquisa baseia-se no uso de técnicas de processamento de linguagem natural para estruturar uma base analítica e treinar modelos com representações numéricas dos textos.

Nesta entrega, o grupo define a linguagem de programação empregada (Python), detalha como a base foi analisada exploratoriamente (EDA), descreve o tratamento previsto para extração, limpeza, tokenização e segmentação (chunking) para preparação do conjunto de treino e avaliação. Adicionalmente, são apresentadas as bases teóricas dos métodos utilizados (TF-IDF, Naive Bayes Multinomial e Regressão Logística multiclasse) e a forma de cálculo da métrica principal de avaliação (acurácia), incluindo complementos para leitura em cenários multiclasse.

% --------------------------------------------------------------------
\clearpage
\section{Definição da Linguagem de Programação}

\subsection{Linguagem Escolhida}

A linguagem de programação escolhida para o desenvolvimento do projeto é o \textbf{Python}.

\subsection{Justificativa da Escolha}

A decisão por Python se justifica por sua adequação ao escopo do projeto, que envolve manipulação de dados e processamento de texto. Em termos práticos, Python permite:

\begin{enumerate}[leftmargin=1.2cm]
    \item Construir um pipeline reprodutível com leitura/escrita de artefatos (por exemplo: PDFs, textos extraídos, chunks e CSVs consolidados).
    \item Aplicar rotinas de extração textual a partir de PDFs, quando necessário, por meio de bibliotecas específicas.
    \item Representar e transformar textos em vetores numéricos por meio de \textit{TF-IDF} e treinar classificadores supervisionados com \texttt{scikit-learn}.
    \item Avaliar modelos com métricas adequadas para classificação multiclasse, como acurácia e métricas complementares (precision, recall, F1 e matriz de confusão).
    \item Organizar a implementação em módulos claros (extração, pré-processamento, modelagem e avaliação), facilitando a colaboração do grupo e a rastreabilidade das decisões.
\end{enumerate}

\subsection{Bibliotecas Previstas}

As bibliotecas usadas/previstas no desenvolvimento do pipeline incluem:

\begin{table}[h]
\centering
\renewcommand{\arraystretch}{1.25}
\begin{tabularx}{\textwidth}{>{\bfseries}p{3.6cm} Y}
\toprule
Biblioteca & Finalidade no projeto \\
\midrule
\texttt{pandas} & Organização tabular da base, leitura de CSVs e estruturação de metadados \\
\texttt{numpy} & Operações numéricas auxiliares \\
\texttt{matplotlib}/\texttt{seaborn} & Visualizações para EDA \\
\texttt{scikit-learn} & TF-IDF, modelos (Naive Bayes, Regressão Logística) e avaliação \\
\texttt{requests} & Coleta automatizada (etapa de coleta, quando aplicável) \\
\texttt{pdfplumber} & Extração de texto de PDFs \\
\bottomrule
\end{tabularx}
\caption{Bibliotecas utilizadas/previstas para o projeto}
\label{tab:bibliotecas}
\end{table}

% --------------------------------------------------------------------
\clearpage
\section{Análise Exploratória da Base de Dados}

\subsection{Base Escolhida}

A base principal do projeto é composta por boletins epidemiológicos em \textbf{PDF} relacionados a arboviroses urbanas (incluindo Dengue, Zika e Chikungunya). Para esta etapa, os boletins utilizados foram extraídos do seguinte endereço oficial:

\vspace{0.2cm}
\noindent
\textbf{Fonte oficial dos boletins (COVISA/SMS-SP):} \\
\url{https://prefeitura.sp.gov.br/web/saude/w/vigilancia_em_saude/boletim_covisa/arboviroses}

\vspace{0.2cm}
\noindent
Observação do próprio site: recomenda-se utilizar o \textit{último boletim publicado}, e boletins anteriores devem ser entendidos como uma fotografia do momento em que foram publicados (não é recomendada comparação direta entre anos diferentes com o boletim vigente).

\subsection{Objetivo da Análise Exploratória}

A análise exploratória tem como objetivo compreender a estrutura do corpus textual, identificar padrões relevantes para modelagem e orientar a preparação dos dados (limpeza, tokenização e segmentação). Além da descrição, a EDA serve para verificar:

\begin{enumerate}[leftmargin=1.2cm]
    \item Quantidade de documentos disponíveis e distribuição de tamanho.
    \item Distribuição de vocabulário e termos predominantes.
    \item Necessidade de segmentação do texto (evitar mistura temática em uma única amostra).
    \item Presença de ruídos típicos da extração de PDFs e como mitigá-los.
    \item Parâmetros de chunking adequados para viabilizar treino e avaliação supervisionada.
\end{enumerate}

\subsection{Unidade de Análise}

A unidade de análise primária é o documento (boletim completo). Entretanto, para aumentar a capacidade do modelo em capturar temas locais e reduzir mistura temática, a etapa prevê a segmentação do texto em \textbf{chunks} (trechos) a partir de janelas de tokens.

\subsection{Resultados Reais da EDA (base processada)}

Com os boletins disponíveis, foram observados os seguintes resultados após extração, limpeza e preparação do texto (conforme artefatos em \texttt{data/processed/}):

\begin{itemize}[leftmargin=1.2cm]
    \item \textbf{Quantidade de PDFs:} 11
    \item \textbf{Chunks gerados:} 60
    \item \textbf{Parâmetros de chunking:} \texttt{chunk\_size=500}, \texttt{stride=200}, \texttt{min\_fill\_ratio=0.6}
    \item \textbf{Tokens por chunk:} média 487.62, mínimo 356, máximo 500
    \item \textbf{Vocabulário (termos únicos após filtragem):} 440
\end{itemize}

\subsubsection{Termos predominantes}

Os termos mais frequentes do corpus, após filtragem (remoção de números puros e tokens muito curtos), reforçam que o conteúdo é dominado por vocabulário epidemiológico e pela estrutura recorrente dos boletins:

\begin{table}[h]
\centering
\renewcommand{\arraystretch}{1.15}
\begin{tabularx}{\textwidth}{>{\bfseries}l Y}
\toprule
Termo & Ocorrências \\
\midrule
casos & 799 \\
sao & 669 \\
ate & 600 \\
confirmados & 569 \\
sinan & 550 \\
municipio & 509 \\
inicio & 507 \\
quadro & 502 \\
online & 494 \\
ano & 485 \\
sintomas & 485 \\
entre & 476 \\
vila & 461 \\
paulo & 436 \\
segundo & 432 \\
\bottomrule
\end{tabularx}
\caption{Top termos mais frequentes (EDA)}
\label{tab:toptermos}
\end{table}

\subsubsection{Implicações para a modelagem}

Com base na EDA, o grupo identifica as seguintes implicações práticas:

\begin{enumerate}[leftmargin=1.2cm]
    \item A recorrência de termos sugere que a representação por \textbf{TF-IDF} é adequada para capturar termos discriminantes entre categorias.
    \item A presença de conteúdo relacionado a mais de uma arbovirose no mesmo boletim aumenta a chance de mistura temática; por isso, a segmentação em chunks favorece a qualidade das amostras.
    \item A similaridade entre boletins (formatos e seções recorrentes) motiva uma estratégia de rotulagem por trechos (chunks), reduzindo mistura de temas em uma única amostra.
\end{enumerate}

\subsubsection{Ruído da extração e limpeza aplicada}

Como PDFs podem introduzir artefatos na extração textual, foram aplicadas rotinas de limpeza antes do chunking:

\begin{itemize}[leftmargin=1.2cm]
    \item Uniao de hifens de quebra de linha (ex.: \texttt{transmissao-continua} $\rightarrow$ \texttt{transmissaocontinua}).
    \item Remoção de linhas que se tornam apenas números (ex.: numeração de páginas/rodapés).
    \item Normalização de \textit{whitespace} (remoção de múltiplos espaços e excesso de quebras de linha).
    \item Normalização de acentos para ASCII (portabilidade dos artefatos no repositório).
    \item Tokenização por \textit{regex} e filtragem (remoção de números puros e tokens com menos de 3 caracteres).
\end{itemize}

% --------------------------------------------------------------------
\clearpage
\section{Tratamento da Base de Dados}

\subsection{Visão Geral do Pipeline}

O tratamento da base de dados compreende um pipeline composto por:

\begin{enumerate}[leftmargin=1.2cm]
    \item Coleta/importação dos PDFs para a estrutura do repositório.
    \item Extração textual a partir dos PDFs.
    \item Limpeza e normalização do texto.
    \item Tokenização e filtragem para compatibilidade com TF-IDF.
    \item Segmentação (chunking) em janelas de tokens.
    \item Criação de um dataset consolidado (CSV) para modelagem.
    \item Rotulagem supervisionada (quando as categorias forem definidas pelo grupo).
    \item Vetorização e treino dos modelos.
    \item Avaliação das predições no conjunto de teste.
\end{enumerate}

\subsection{Preparação (extração, limpeza, tokenização e chunking)}

O processo foi implementado no módulo \texttt{src/preprocessing.py} e gera os seguintes artefatos:

\begin{itemize}[leftmargin=1.2cm]
    \item \texttt{data/processed/extracted\_text/}: texto limpo por documento (\texttt{.md}).
    \item \texttt{data/processed/chunks/}: texto por chunk (\texttt{.md}).
    \item \texttt{data/processed/chunks\_dataset.csv}: dataset consolidado para modelagem.
    \item \texttt{data/processed/eda\_summary.*}: resumo estatístico da EDA (tokens, vocabulário e top termos).
\end{itemize}

\subsubsection{Chunking e identificação de amostras}

Cada chunk recebe um identificador no formato:

\[
 \text{chunk\_id} = \text{doc\_id}\_\text{chunk\_index}
 \]

Com:
\begin{itemize}[leftmargin=1.2cm]
    \item \texttt{doc\_id}: derivado do nome do arquivo PDF (sem extensão).
    \item \texttt{chunk\_index}: índice sequencial do chunk dentro do documento.
\end{itemize}

\subsubsection{Dataset final para modelagem}

O arquivo \texttt{data/processed/chunks\_dataset.csv} contém (no mínimo) as colunas:

\begin{itemize}[leftmargin=1.2cm]
    \item \texttt{chunk\_id}
    \item \texttt{doc\_id}
    \item \texttt{chunk\_index}
    \item \texttt{token\_count}
    \item \texttt{chunk\_text}
    \item \texttt{text\_md\_path}
\end{itemize}

\subsection{Rotulagem supervisionada (labels)}

Para que o treino supervisionado ocorra, será necessário criar:

\begin{itemize}[leftmargin=1.2cm]
    \item \texttt{data/processed/labels.csv} com colunas:
    \begin{itemize}
        \item \texttt{chunk\_id}
        \item \texttt{label}
    \end{itemize}
\end{itemize}

\subsection{Vetorização e Treinamento}

Com o dataset rotulado, os textos serão representados numericamente por TF-IDF e treinados em dois classificadores:

\begin{enumerate}[leftmargin=1.2cm]
    \item \textbf{Baseline (referência inicial)}: \texttt{MultinomialNB} (\textit{Naive Bayes Multinomial}).
    \item \textbf{Modelo principal}: \texttt{LogisticRegression} em modo multiclasse (\textit{multinomial}) para classificação discriminativa.
\end{enumerate}

\subsubsection{Parâmetros de TF-IDF}

O \texttt{TfidfVectorizer} utiliza:

\begin{itemize}[leftmargin=1.2cm]
    \item \texttt{ngram\_range=(1,2)} (unigramas e bigramas)
    \item \texttt{min\_df=2} e \texttt{max\_df=0.95}
    \item \texttt{max\_features=20000}
    \item \texttt{stop\_words=None} (mantém termos do domínio)
\end{itemize}

\subsubsection{Divisão de dados (treino/validação/teste)}

Quando labels estiverem disponíveis, a divisão será estratificada por classe:

\begin{itemize}[leftmargin=1.2cm]
    \item 70\% treino
    \item 15\% validação
    \item 15\% teste final
\end{itemize}

% --------------------------------------------------------------------
\clearpage
\section{Definição e Descrição das Bases Teóricas dos Métodos}

\subsection{Aprendizado Supervisionado para Classificação Textual}

O problema será formulado como uma tarefa de \textbf{classificação supervisionada}: os textos vetorizados são associados a rótulos (categorias temáticas previamente atribuídas). O objetivo é aprender padrões que permitam generalizar a classificação para novos textos.

\subsection{Base Teórica da Representação TF-IDF}

A técnica TF-IDF atribui peso maior a termos que aparecem com frequência no documento, mas que não são excessivamente comuns no corpus. Formalmente:

\[
\operatorname{TF\mbox{-}IDF}(t,d) = \operatorname{tf}(t,d)\cdot \log\left(\frac{N}{\operatorname{df}(t)}\right)
\]

em que:
\begin{itemize}[leftmargin=1.2cm]
    \item $N$ é o número total de documentos no corpus;
    \item $\operatorname{df}(t)$ é o número de documentos em que o termo $t$ aparece;
    \item $\operatorname{tf}(t,d)$ é a frequência do termo $t$ no documento $d$.
\end{itemize}

\subsection{Modelo 1: Naive Bayes Multinomial}

O \textit{Naive Bayes Multinomial} é um classificador probabilístico baseado na probabilidade de uma classe dado o conjunto de palavras observadas no texto. Sua decisão é expressa como:

\[
\hat{c} = \arg\max_{c \in C} P(c)\prod_{i=1}^{n}P(w_i \mid c)
\]

\subsection{Modelo 2: Regressão Logística Multiclasse}

A \textit{Regressão Logística Multiclasse} estima probabilidades por classe. Com função softmax:

\[
P(y = k \mid x) = \frac{e^{\beta_k^\top x}}{\sum_{j=1}^{K} e^{\beta_j^\top x}}
\]

\subsection{Estratégia Analítica Adotada}

O grupo utiliza o \textit{Naive Bayes Multinomial} como modelo de referência e a \textit{Regressão Logística Multiclasse} como modelo principal para comparação.

% --------------------------------------------------------------------
\clearpage
\section{Definição e Descrição de Como Será Calculada a Acurácia}

\subsection{Métrica Principal}

A métrica principal de avaliação será a \textbf{acurácia}, definida como a proporção de previsões corretas em relação ao total de previsões:

\[
\text{Acurácia} = \frac{\text{número de classificações corretas}}{\text{número total de classificações}}
\]

Essa métrica será calculada no conjunto de teste (com rótulos reais) após o treino e a seleção usando apenas a estratégia de validação, evitando vazamento de informação.

\subsection{Métricas Complementares}

Além da acurácia, o grupo calculará métricas complementares para classificação multiclasse:
\begin{itemize}[leftmargin=1.2cm]
    \item \textbf{Precision} por classe e média macro;
    \item \textbf{Recall} por classe e média macro;
    \item \textbf{F1-score} por classe e média macro;
    \item \textbf{Matriz de confusão} para análise de erros entre categorias.
\end{itemize}

% --------------------------------------------------------------------
\clearpage
\section{Considerações Finais}

Com base na EDA realizada sobre boletins epidemiológicos em PDF (11 documentos, segmentados em 60 chunks), o grupo definiu uma abordagem metodológica coerente para classificação textual supervisionada. O pipeline de tratamento (extração, limpeza, tokenização, chunking e construção do dataset) fornece uma base reprodutível para o treino de modelos via TF-IDF.

\vspace{0.3cm}
\noindent
\textbf{Fonte oficial dos boletins:} \\
\url{https://prefeitura.sp.gov.br/web/saude/w/vigilancia_em_saude/boletim_covisa/arboviroses}

